\SourceLeaf{536}senkrechten Spalte jeder einzelnen Tabelle die drei niedrigsten Ziffern des Quadrates und in dem Innern die vier oder fünf höheren Ziffern des Quadrates.
\par Die Tabelle ‘\textit{Indices der Primzahlen im Höhern Zahlenreiche}’ enthält in der obersten horizontalen Reihe den jedesmaligen Modulus, in der zweiten Reihe die zur Anwendung gekommene Basis, in der ersten senkrechten Spalte die Restzahlen und im Innern die Indices. Die Sterne $*$ bezeichnen die Reste Null.
\par Die Handschrift des Bruchstückes der ‘\textit{Sectio octava. Quarundam disquisitionum ad circuli sectionem pertinentium uberior consideratio}’ scheint der Zeit der Umarbeitung der ‘Analysis residuorum’ in die ‘Disquisitiones arithmeticae’ anzugehören. Die Briefe von \textsc{Gauss} an \textsc{Dirichlet} bestätigen die Ansicht, dass \textsc{Gauss} auch in der höheren Arithmetik erheblich mehr entdeckt hat als im Nachlasse sich findet.
\begin{flushright}\textsc{Schering}.\end{flushright}
\BookRule{.25\textwidth}
